\documentclass[12pt,a4paper]{article}
\usepackage[a4paper, left=20mm, right=22mm, top=32mm, bottom=22mm, headheight=28mm, headsep=6mm, footskip=12mm]{geometry}
\usepackage{fontspec}
\IfFontExistsTF{Times New Roman}{\setmainfont{Times New Roman}}{\setmainfont{Liberation Serif}}
\usepackage{graphicx}
\usepackage{fancyhdr}
\usepackage{tabularx}
\usepackage{booktabs}
\usepackage{array}
\usepackage{xcolor}

% Background watermark centered on every page
\setlength{\unitlength}{1pt}
\AddToHook{shipout/background}{%
  \put(65,-653){%
    \includegraphics[width=466pt,height=465pt]{assets/watermark.png}%
  }%
}

% Header and Footer configuration
\pagestyle{fancy}
\fancyhf{}
\fancyhead[L]{%
  \raisebox{-0.35\height}{\includegraphics[width=22mm,height=22mm,keepaspectratio]{assets/header_logo.png}}%
  \hspace{4mm}%
  \parbox[b]{120mm}{%
    {\fontsize{11.5pt}{14pt}\selectfont \textbf{School of Electrical and Computer Engineering}}\par\vspace{2pt}%
    {\fontsize{11.5pt}{14pt}\selectfont University of Tehran}%
  }%
}
\fancyfoot[R]{\fontsize{11pt}{13pt}\selectfont Page | \thepage}
\renewcommand{\headrulewidth}{0pt}
\renewcommand{\footrulewidth}{0pt}

\setlength{\parindent}{0pt}
\setlength{\parskip}{4pt}
\linespread{1.15}

\begin{document}

% ================= PAGE 1: COVER =================
\thispagestyle{empty}
\begin{center}
  \vspace*{10mm}
  \includegraphics[width=85mm,height=85mm,keepaspectratio]{assets/cover_logo.png}\par\vspace{12mm}
  {\fontsize{24pt}{28pt}\selectfont \textbf{University of Tehran}}\par\vspace{8mm}
  {\fontsize{18pt}{22pt}\selectfont \textbf{Tehran}}\par\vspace{14mm}
  {\fontsize{18pt}{22pt}\selectfont \textbf{School of Electrical and Computer Engineering}}\par\vspace{8mm}
  {\fontsize{18pt}{22pt}\selectfont \textbf{Course Description}}\par\vspace{8mm}
  {\fontsize{18pt}{22pt}\selectfont \textbf{\textit{Computer Engineering}}}\par\vspace{8mm}
  {\fontsize{18pt}{22pt}\selectfont \textbf{Bachelor’s Degree}}
\end{center}
\newpage

% ================= PAGE 2: REMARKS =================
\setcounter{page}{1}
\textbf{Remarks:}\par\vspace{4pt}
According to the Academic Instruction of the University, in the academic system, each Theoretical credit is presented in 16 hours, each Practical credit in 32 hours, and each Workshop credit in 48 hours.\par\vspace{8pt}
There are some different types of courses passed by a bachelor/master/Ph.D. student at this university, as listed below:
\begin{enumerate}
  \setlength{\itemsep}{1pt}
  \setlength{\parskip}{0pt}
  \item Basic
  \item Major
  \item Specialized
  \item Elective
  \item General
  \item Workshop
  \item Internship
  \item Project
\end{enumerate}
\vspace{4pt}
Each of these types contains some courses to be passed according to the educational regulations passed by the Ministry of Science, Research and Technology or the Ministry of Health \& Medical Education of Iran.\par\vspace{8pt}
The grading system in this university is from 0 to 20. The minimum passing grade for a course leading to an Associate's Degree or a Bachelor's Degree is 10, for a course leading to a Master's Degree and Medical Sciences is 12, and for a course leading to a Ph.D. Degree is 14.
\newpage

% ================= CURRICULUM OVERVIEW =================
\textbf{\LARGE Curriculum Overview: Semester-by-Semester Study Plan}\par\vspace{6pt}
The following study plan outlines the curriculum courses undertaken across all academic semesters, including course categories and credit allocations. This document provides official course descriptions and syllabi for all included courses.\par\vspace{10pt}
\noindent{\large\textbf{Semester 1: Academic Year 2023–2024 (1st Semester)}}\par\vspace{4pt}
\noindent\begin{tabularx}{\textwidth}{@{} X l r @{}}
\toprule
\textbf{Course Title} & \textbf{Course Type} & \textbf{Credits} \\
\midrule
Introduction to Computing Systems and Programming & Specialized | Theoretical & 4 \\
Physics 1 & Basic | Theoretical & 3 \\
Calculus 1 & Basic | Theoretical & 3 \\
Islamic Thought 1 (Origins and Resurrection) & General | Theoretical & 2 \\
Analytical history of beginning Islam & General | Theoretical & 2 \\
Persian Language & General | Theoretical & 3 \\
\midrule
\multicolumn{2}{@{}l}{\textbf{Total Semester Credits: 17}} & \textbf{Cumulative: 17} \\
\bottomrule
\end{tabularx}\par\vspace{14pt}
\noindent{\large\textbf{Semester 2: Academic Year 2023–2024 (2nd Semester)}}\par\vspace{4pt}
\noindent\begin{tabularx}{\textwidth}{@{} X l r @{}}
\toprule
\textbf{Course Title} & \textbf{Course Type} & \textbf{Credits} \\
\midrule
Advanced Programming & Specialized | Theoretical & 3 \\
Differential Equations & Basic | Theoretical & 3 \\
Physics 2 & Basic | Theoretical & 3 \\
Calculus 2 & Basic | Theoretical & 3 \\
Discrete Mathematics & Major | Theoretical & 3 \\
Islamic Thought 2 (Prophethood and Imamat) & General | Theoretical & 2 \\
\midrule
\multicolumn{2}{@{}l}{\textbf{Total Semester Credits: 17}} & \textbf{Cumulative: 34} \\
\bottomrule
\end{tabularx}\par\vspace{14pt}
\noindent{\large\textbf{Semester 3: Academic Year 2024–2025 (1st Semester)}}\par\vspace{4pt}
\noindent\begin{tabularx}{\textwidth}{@{} X l r @{}}
\toprule
\textbf{Course Title} & \textbf{Course Type} & \textbf{Credits} \\
\midrule
Engineering Probability and Statistics & Basic | Theoretical & 3 \\
Logic Circuits & Specialized | Theoretical & 3 \\
Electric Circuits & Specialized | Theoretical & 3 \\
Islamic Revolution in Iran & General | Theoretical & 2 \\
Family Schematization and Population & General | Theoretical & 2 \\
Data Structures and Algorithms & Major | Theoretical & 3 \\
\midrule
\multicolumn{2}{@{}l}{\textbf{Total Semester Credits: 16}} & \textbf{Cumulative: 50} \\
\bottomrule
\end{tabularx}\par\vspace{14pt}
\noindent{\large\textbf{Semester 4: Academic Year 2024–2025 (2nd Semester)}}\par\vspace{4pt}
\noindent\begin{tabularx}{\textwidth}{@{} X l r @{}}
\toprule
\textbf{Course Title} & \textbf{Course Type} & \textbf{Credits} \\
\midrule
Logic Circuits Laboratory & Specialized | Practical & 1 \\
Engineering Mathematics & Specialized | Theoretical & 3 \\
Computer Workshop & Basic | Practical & 1 \\
Artificial Intelligence & Specialized | Theoretical & 3 \\
Computer Architecture & Specialized | Theoretical & 3 \\
Machines and Language Theory & Specialized | Theoretical & 3 \\
Electrical Measurment and Circuit Laboratory & Specialized | Practical & 1 \\
Foundations of Data Science & Elective | Theoretical & 3 \\
\midrule
\multicolumn{2}{@{}l}{\textbf{Total Semester Credits: 18}} & \textbf{Cumulative: 68} \\
\bottomrule
\end{tabularx}\par\vspace{14pt}
\noindent{\large\textbf{Semester 5: Academic Year 2025–2026 (1st Semester)}}\par\vspace{4pt}
\noindent\begin{tabularx}{\textwidth}{@{} X l r @{}}
\toprule
\textbf{Course Title} & \textbf{Course Type} & \textbf{Credits} \\
\midrule
Operating Systems Laboratory & Specialized | Practical & 1 \\
Operating System & Specialized | Theoretical & 3 \\
Computer Aided Digital System & Specialized | Theoretical & 3 \\
Signals and Systems & Specialized | Theoretical & 3 \\
English Language & General | Theoretical & 3 \\
Algorithm Design & Major | Theoretical & 3 \\
Basucs of Machine Learning & Elective | Theoretical & 3 \\
\midrule
\multicolumn{2}{@{}l}{\textbf{Total Semester Credits: 19}} & \textbf{Cumulative: 87} \\
\bottomrule
\end{tabularx}\par\vspace{14pt}
\noindent{\large\textbf{Semester 6: Academic Year 2025–2026 (2nd Semester)}}\par\vspace{4pt}
\noindent\begin{tabularx}{\textwidth}{@{} X l r @{}}
\toprule
\textbf{Course Title} & \textbf{Course Type} & \textbf{Credits} \\
\midrule
Systems Analysis and Design & Elective | Theoretical & 3 \\
Technical English & Specialized | Theoretical & 2 \\
Computer Networks & Specialized | Theoretical & 3 \\
Database Design & Elective | Theoretical & 3 \\
Principles of cognitive science \textit{\textcolor{red!75!black}{(Not Included in Syllabus)}} & Specialized | Theoretical & 3 \\
Physical Education & General | Practical & 1 \\
Compiler design and programming languages & Specialized | Theoretical & 3 \\
\midrule
\multicolumn{2}{@{}l}{\textbf{Total Semester Credits: 18}} & \textbf{Cumulative: 105} \\
\bottomrule
\end{tabularx}\par\vspace{14pt}
\noindent{\large\textbf{Summer Semester (2025–2026)}}\par\vspace{4pt}
\noindent\begin{tabularx}{\textwidth}{@{} X l r @{}}
\toprule
\textbf{Course Title} & \textbf{Course Type} & \textbf{Credits} \\
\midrule
Industrial Training & Training | Practical & 3 \\
\midrule
\multicolumn{2}{@{}l}{\textbf{Total Semester Credits: 0}} & \textbf{Cumulative: 105} \\
\bottomrule
\end{tabularx}\par\vspace{14pt}
\noindent{\large\textbf{Semester 7: Academic Year 2026–2027 (1st Semester)}}\par\vspace{4pt}
\noindent\begin{tabularx}{\textwidth}{@{} X l r @{}}
\toprule
\textbf{Course Title} & \textbf{Course Type} & \textbf{Credits} \\
\midrule
Software Engineering & Elective | Theoretical & 3 \\
Database Laboratory & Elective | Practical & 1 \\
Information Technology Project Management & Specialized | Theoretical & 3 \\
Computer Networks Security & Elective | Theoretical & 3 \\
Introduction to Software Testing & Elective | Theoretical & 3 \\
Thematic interpretation of the Quran & General | Theoretical & 2 \\
Sport 1 & General | Practical & 1 \\
Generative Models \textit{\textcolor{red!75!black}{(Not Included in Syllabus)}} & Specialized | Theoretical & 3 \\
\midrule
\multicolumn{2}{@{}l}{\textbf{Total Semester Credits: 19}} & \textbf{Cumulative: 124} \\
\bottomrule
\end{tabularx}\par\vspace{14pt}
\vspace{3pt}\noindent\textit{\textcolor{red!75!black}{* Courses marked (Not Included in Syllabus) are present in the academic transcript but not currently archived in the course catalogue.}}\par
\newpage

% ================= COURSES BODY =================
\newpage
\vspace*{\fill}
\begin{center}
  {\fontsize{28pt}{34pt}\selectfont \textbf{SEMESTER 1}}\par\vspace{12pt}
  {\fontsize{13pt}{17pt}\selectfont Academic Year 2023–2024 (1st Semester)}
\end{center}
\vspace*{\fill}

\newpage
\textbf{Title of the Course:} Introduction to Computing Systems and Programming\\
\textbf{Number of Credits:} 4\\
\textbf{Type of the Course:} Specialized | Theoretical\\
\textbf{Training hours:} 64h Theoretical\par\vspace{4pt}
\textbf{Course Objectives:}\par\vspace{2pt}
Explanation of main computer parts, the concept of software and hardware, algorithms design, and
an introduction to a structured computer programming language.\par\vspace{4pt}
\textbf{Syllabus of Courses:}\par\vspace{2pt}
Computer fundamentals, including the role of computers in the modern world and their\par
applications. Introduction to main computer parts and the concept of software and hardware.\par
Native data types (Integrals and Floating points). Algorithm design and problem-solving principles\par
(Flow, selection, and loops). Implanting the algorithm by Pseudo Code. An introduction to a\par
structured programming language. Constants, variables, computational and logical expressions.\par
language-specific instructions. Conditional structures. Loops, Arrays, vectors and matrices,\par
functions, and procedures. I/O and files operations. Implementing simple algorithms by the\par
programming language. Searching and sorting. Advanced program design.\par
Programming language: C/C++\par\vspace{4pt}
\textbf{TextBooks:}\par\vspace{2pt}
1. Bjarne Stroustrup, "The C++ Programming Language", Addison Wesley, 2004.\par
2. Stanley B Lippman, "The C++ Primer", Addison Wesley, 2005.\par
3. Paul Deitel \& Harvey Deitel, “C++ How to Program (Early Objects Version)”, Prentice Hall; 9\par
edition, 2013.\par
\vspace{4pt}

\newpage
\textbf{Title of the Course:} Physics 1\\
\textbf{Number of Credits:} 3\\
\textbf{Type of the Course:} Basic | Theoretical\\
\textbf{Training hours:} 48h Theoretical\par\vspace{4pt}
\textbf{Course Objectives:}\par\vspace{2pt}
The main goal of fundamental courses in physics is to provide tools by which students can learn
how to read scientific material effectively, identify fundamental concepts, reason through scientific
questions, and solve quantitative problems. Physics 1 is the first course of this set. This course
covers the fundamental concepts in Classical Mechanics and Thermodynamics.\par\vspace{4pt}
\textbf{Syllabus of Courses:}\par\vspace{2pt}
\vspace{2pt}\textbf{- Units, Physical Quantities, and Vectors:}\par
The Nature of Physics, Solving Physics Problems, Standards and Units, Unit Consistency and\par
Conversions, Uncertainty and Significant Figures, Estimates and Orders of Magnitude, Vectors\par
and Vector Addition, Components of Vectors, Unit Vectors, Products of Vectors\par
\vspace{2pt}\textbf{- Motion along a Straight Line:}\par
Displacement, Time, and Average Velocity, Instantaneous Velocity, Average and Instantaneous\par
Acceleration, Motion with Constant Acceleration, Freely Falling Bodies, Velocity and Position by\par
Integration\par
\vspace{2pt}\textbf{- Motion in Two or Three Dimensions:}\par
Position and Velocity Vectors, The Acceleration Vector, Projectile Motion, Motion in a Circle,\par
Relative Velocity\par
\vspace{2pt}\textbf{- Newton’s Laws of Motion:}\par
Force and Interactions, Newton’s First Law, Newton’s Second Law, Mass and Weight, Newton’s\par
Third Law, Free-Body Diagrams\par
\vspace{2pt}\textbf{- Applying Newton’s Laws:}\par
Using Newton’s First Law: Particles in Equilibrium, Using Newton’s Second Law: Dynamics of\par
Particles, Frictional Forces, Dynamics of Circular Motion, The Fundamental Forces of Nature\par
\vspace{2pt}\textbf{- Work and Kinetic Energy:}\par
Work, Kinetic Energy, and the Work-Energy Theorem, Work and Energy with Varying Forces,\par
Power\par
\vspace{2pt}\textbf{- Potential Energy and Energy Conservation:}\par
Gravitational Potential Energy, Elastic Potential Energy, Conservative and Non-conservative\par
Forces, Force and Potential Energy, Energy Diagrams\par
\vspace{2pt}\textbf{- Momentum, Impulse, and Collisions:}\par
Momentum and Impulse, Conservation of Momentum, Momentum Conservation and Collisions,\par
Elastic Collisions, Center of Mass, Rocket Propulsion\par
\vspace{2pt}\textbf{- Rotation of Rigid Bodies:}\par
Angular Velocity and Acceleration, Rotation with Constant Angular Acceleration, Relating Linear\par
and Angular Kinematics, Energy in Rotational Motion, Parallel-Axis Theorem, Moment-of-Inertia\par
Calculations\par
\vspace{2pt}\textbf{- Dynamics of Rotational Motion:}\par
Torque, Torque and Angular Acceleration for a Rigid-Body, Rigid-Body Rotation About a Moving\par
Axis, Work and Power in Rotational Motion, Angular Momentum, Conservation of Angular\par
Momentum, Gyroscopes, and Precession\par
\vspace{2pt}\textbf{- Equilibrium and Elasticity:}\par
Conditions for Equilibrium, Center of Gravity, Solving Rigid-Body Equilibrium Problems, Stress,\par
Strain, Elastic Moduli, Elasticity, and Plasticity\par
\vspace{2pt}\textbf{- Fluid Mechanics:}\par
Density, Pressure in a Fluid, Buoyancy, Fluid Flow, Bernoulli’s Equation, Viscosity, and\par
Turbulence\par
\vspace{2pt}\textbf{- Gravitation:}\par
Newton’s Law of Gravitation, Weight, Gravitational Potential Energy, The Motion of Satellites,\par
Kepler’s Laws and the Motion of Planets, Spherical Mass Distributions, Apparent Weight and the\par
Earth’s Rotation, Black Holes\par
\vspace{2pt}\textbf{- Periodic Motion:}\par
Describing Oscillation, Simple Harmonic Motion, Energy in Simple Harmonic Motion,\par
Applications of Simple Harmonic Motion, The Simple Pendulum, The Physical Pendulum,\par
Damped Oscillations, Forced Oscillations, and Resonance\par
\vspace{2pt}\textbf{- Temperature and Heat:}\par
Temperature and Thermal Equilibrium, Thermometers and Temperature Scales, Gas\par
Thermometers and the Kelvin Scale, Thermal Expansion, Quantity of Heat, Calorimetry and Phase\par
Changes, Mechanisms of Heat Transfer\par
\vspace{2pt}\textbf{- Thermal Properties of Matter:}\par
Equations of State, Molecular Properties of Matter, Kinetic-Molecular Model of an Ideal Gas,\par
Heat Capacities, Molecular Speeds, Phases of Matter\par
\vspace{2pt}\textbf{- The First Law of Thermodynamics:}\par
Thermodynamic Systems, Work Done During Volume Changes, Paths Between Thermodynamic\par
States, Internal Energy and the First Law of Thermodynamics, Kinds of Thermodynamic\par
Processes, Internal Energy of an Ideal Gas, Heat Capacities of an Ideal Gas, Adiabatic Processes\par
for an Ideal Gas\par
\vspace{2pt}\textbf{- The Second Law of Thermodynamics:}\par
Directions of Thermodynamic Processes, Heat Engines, Internal-Combustion Engines,\par
Refrigerators, The Second Law of Thermodynamics, The Carnot Cycle, Entropy, Microscopic\par
Interpretation of Entropy\par\vspace{4pt}
\textbf{TextBooks:}\par\vspace{2pt}
1. Fundamentals of Physics, David Halliday, Jearl Walker, Robert Resnick, 10th edition\par
\vspace{4pt}

\newpage
\textbf{Title of the Course:} Calculus 1\\
\textbf{Number of Credits:} 3\\
\textbf{Type of the Course:} Basic | Theoretical\\
\textbf{Training hours:} 48h Theoretical\par\vspace{4pt}
\textbf{Course Objectives:}\par\vspace{2pt}
Calculus 1, consisted principally of one-variable Calculus, Functions, Derivatives, Integrals,
Integration Methods, Complex Numbers, and Infinite Series.\par\vspace{4pt}
\textbf{Syllabus of Courses:}\par\vspace{2pt}
Cartesian coordinates; polar coordinates; complex numbers; addition, product, root \& geometrical\par
representation of complex numbers; polar representation of complex numbers; function; functions\par
algebra; limit and relevant theorems; infinite limit and limit in infinite; left-hand and right-hand\par
limit; connectivity; derivative; derivation formula; inverse function and its derivative;\par
trigonometric functions derivative and their inverse functions; Rolle’s theorem; mean theorem;\par
Taylor expansion; geometrical and physical applications of derivative; curves and acceleration in\par
polar coordinates; application of derivative in approximation of equations roots; definition of\par
integral of continuous functions and piecewise continuous; basic theorems of differential \&\par
integral arithmetic; primitive function; approximate methods of integral estimate; application of\par
integral in computation of area, length of curve, moment, center of gravity and labor ….(in\par
Cartesian and polar coordinates); logarithm and exponential function and their derivative;\par
hyperbolic functions; integration methods such as change of variable, component and\par
decomposition of fractions; transform of special variables of sequence and numerical series and\par
relevant theorems; power series and Taylor theorem with remainder.\par\vspace{4pt}
\textbf{TextBooks:}\par\vspace{2pt}
1. Tom M. Apostol, “Calculus, Vol. 1: One-Variable Calculus, with an Introduction to Linear\par
Algebra”, John Wiley, 2012.\par
2. George B. Thomas Jr., Maurice D. Weir, Joel Hass, "Thomas' Calculus (12th Edition)", 2009\par
\vspace{4pt}

\newpage
\nopagebreak
\textbf{Title of the Course:} Islamic Thoughts 1 (Beginning and Resurrection)\\
\textbf{Number of Credits:} 2\\
\textbf{Type of the Course:} General | Theoretical\\
\textbf{Training hours:} 32h Theoretical\par\vspace{4pt}
\textbf{Course Objectives:}\par\vspace{2pt}
Explain the necessity of Muslim students to pay attention to religion and religious categories,
deepen and expand information and strengthen students' theism and faith in the field of theology
and resurrection.\par\vspace{4pt}

\par\vspace{14pt}
\nopagebreak
\textbf{Title of the Course:} Analytical History of Beginning Islam\\
\textbf{Number of Credits:} 2\\
\textbf{Type of the Course:} General | Theoretical\\
\textbf{Training hours:} 32h Theoretical\par\vspace{4pt}
\textbf{Course Objectives:}\par\vspace{2pt}
Familiarity of students with how Islam emerged and spread and the way of transition from an
ignorant society to an Islamic society.\par\vspace{4pt}

\par\vspace{14pt}
\nopagebreak
\textbf{Title of the Course:} Persian Language\\
\textbf{Number of Credits:} 3\\
\textbf{Type of the Course:} General | Theoretical\\
\textbf{Training hours:} 48h Theoretical\par\vspace{4pt}
\textbf{Course Objectives:}\par\vspace{2pt}
Familiarity with the Persian language and reading literary texts, poetry, and prose in Iranian
Persian literature.\par\vspace{4pt}

\newpage
\vspace*{\fill}
\begin{center}
  {\fontsize{28pt}{34pt}\selectfont \textbf{SEMESTER 2}}\par\vspace{12pt}
  {\fontsize{13pt}{17pt}\selectfont Academic Year 2023–2024 (2nd Semester)}
\end{center}
\vspace*{\fill}

\newpage
\textbf{Title of the Course:} Advanced Programming\\
\textbf{Number of Credits:} 3\\
\textbf{Type of the Course:} Specialized | Theoretical\\
\textbf{Pre-requisite:} Introduction to Computing Systems and Programming\\
\textbf{Training hours:} 48h Theoretical\par\vspace{4pt}
\textbf{Course Objectives:}\par\vspace{2pt}
The purpose of this course is to study the principles and methods (Polymorphism, Inheritance,
Program Samples, and other paradigms of programming) necessary to produce a computer
program with good-quality specifications. In this regard, after covering the top-down design
method to solve the problem, students will be introduced to the concepts and techniques of objectoriented programming as a tool for managing complexity in medium and large-scale applications.
Program performance, testing, and debugging are focused on throughout the course. An objectoriented programming language (such as C++) is used to teach these principles and methods, and
students are introduced to the concepts and techniques of object-oriented programming and
module-based object design, advanced concepts such as graphical user interface design.\par\vspace{4pt}
\textbf{Syllabus of Courses:}\par\vspace{2pt}
Introduction: Program life cycle, the importance of programming, programming paradigms,\par
Objected orientation\par
Object-oriented programming: Description of Objects, Object recognition, Finding implantable\par
objects in problems, Objects connection, Object-oriented programming language, an Introduction\par
to C++ and its history\par
A review of a functional programming language (C++), Polymorphism, inheritance, program\par
samples\par
Class and Objects, encapsulation and abstract, C++ classes and their parts (private, public, and\par
protected), constructor and destructors, derived classes\par
Polymorphism, polymorphism importance, function polymorphism, operator polymorphism\par
Inheritance: introduction to inheritance, single and multiple inheritances\par
Templates: template importance, function templates, derived data type templates\par
Other programming paradigms: introduction to event-oriented and graphical programming,\par
windows programming, client-side programming\par
Program testing and documentation: the importance of testing and documentation, automatic\par
testing tools, in-line documentation, technical documentation, user manuals\par\vspace{4pt}
\textbf{TextBooks:}\par\vspace{2pt}
1. Bjarne Stroustrup, "The C++ Programming Language", Addison Wesley, 2004.\par
2. Stanley B Lippman, "The C++ Primer", Addison Wesley, 2005.\par
3. Paul Deitel \& Harvey Deitel, “C++ How to Program (Early Objects Version)”, Prentice Hall; 9\par
edition, 2013.\par
\vspace{4pt}

\newpage
\textbf{Title of the Course:} Differential Equations\\
\textbf{Number of Credits:} 3\\
\textbf{Type of the Course:} Basic | Theoretical\\
\textbf{Pre-requisite:} Calculus 1\\
\textbf{Training hours:} 48h Theoretical\par\vspace{4pt}
\textbf{Course Objectives:}\par\vspace{2pt}
In this course, the first and second levels of linear differential equations and some nonlinear
differential equations will be introduced, in addition, students will learn about some numerical and
analytical ways to solve Mathematical Problems.\par\vspace{4pt}
\textbf{Syllabus of Courses:}\par\vspace{2pt}
\vspace{2pt}\textbf{- Introduction}\par
Some Basic Mathematical Models, Direction Fields, Solutions of Some Differential Equations,\par
Classification of Differential Equations, Historical Remarks\par
\vspace{2pt}\textbf{- First Order Differential Equations}\par
Linear Equations with Variable Coefficients, Separable Equations, Modeling with First Order\par
Equations, Differences between Linear and Nonlinear Equations, Autonomous Equations and\par
Population Dynamics, Exact Equations and Integrating Factors, The Existence and Uniqueness\par
Theorem, First Order Difference Equations\par
\vspace{2pt}\textbf{- Second Order Linear Equations}\par
Homogeneous Equations with Constant Coefficients, Fundamental Solutions of Linear\par
Homogeneous Equations, Linear Independence and the Wronskian, Complex Roots of the\par
Characteristic Equation, Repeated Roots, Reduction of Order, Nonhomogeneous Equations,\par
Method of Undetermined Coefficients, Variation of Parameters.\par
\vspace{2pt}\textbf{- Higher Order Linear Equations}\par
General Theory of nth Order Linear Equations, Homogeneous Equations with Constant\par
Coefficients, The Method of Undetermined Coefficients, The Method of Variation of Parameters\par
\vspace{2pt}\textbf{- Series Solutions of Second-Order Linear Equations}\par
Review of Power Series, Series Solutions near an Ordinary Point, Regular Singular Points, Euler\par
Equations, Series Solutions near a Regular Singular Point, Bessel’s Equation\par
\vspace{2pt}\textbf{- The Laplace Transform}\par
Definition of the Laplace Transform, Solution of Initial Value Problems, Step Functions,\par
Differential Equations with Discontinuous Forcing Functions, Impulse Functions, and The\par
Convolution Integral\par
\vspace{2pt}\textbf{- Systems of First-Order Linear Equations}\par
Introduction, Review of Matrices, Systems of Linear Algebraic Equations; Linear Independence,\par
Eigenvalues, Eigenvectors, Basic Theory of Systems of First Order Linear Equations\par
Homogeneous Linear Systems with Constant Coefficients, Complex Eigenvalues, Fundamental\par
Matrices, Repeated Eigenvalues,\par\vspace{4pt}
\textbf{TextBooks:}\par\vspace{2pt}
1. E. Kreyszig, Advanced Engineering Mathematics, 1999\par
2. P. Blanchard, R.L. Devaney, and Co. Hall, Differential Equations, 1st Edition, Brooks/ColePub,\par
1998\par
\vspace{4pt}

\newpage
\textbf{Title of the Course:} Physics 2\\
\textbf{Number of Credits:} 3\\
\textbf{Type of the Course:} Basic | Theoretical\\
\textbf{Pre-requisite:} Physics 1\\
\textbf{Training hours:} 48h Theoretical\par\vspace{4pt}
\textbf{Course Objectives:}\par\vspace{2pt}
The main goal of fundamental courses in physics is to provide tools by which students can learn
how to read scientific material effectively, identify fundamental concepts, reason through scientific
questions, and solve quantitative problems. Physics 2 is the second course of this set. This course
covers the fundamental concepts of Electromagnetism.\par\vspace{4pt}
\textbf{Syllabus of Courses:}\par\vspace{2pt}
\vspace{2pt}\textbf{- Electric Charge and Electric Field:}\par
Electric Charge, Conductors, Insulators, and Induced Charges, Coulomb’s Law, Electric Field and\par
Electric Forces, Electric-Field Calculations, Electric Field Lines, Electric Dipoles\par
\vspace{2pt}\textbf{- Gauss’s Law:}\par
Charge and Electric Flux, Calculating Electric Flux, Gauss’s Law, Applications of Gauss’s Law,\par
Charges on Conductors\par
\vspace{2pt}\textbf{- Electric Potential:}\par
Electric Potential Energy, Electric Potential, Calculating Electric Potential, Equipotential\par
Surfaces, Potential Gradient\par
\vspace{2pt}\textbf{- Capacitance and Dielectrics:}\par
Capacitors and Capacitance, Capacitors in Series and Parallel, Energy Storage in Capacitors and\par
Electric-Field Energy, Dielectrics, Molecular Model of Induced Charge, Gauss’s Law in\par
Dielectrics\par
\vspace{2pt}\textbf{- Current, Resistance, and Electromotive Force:}\par
Current, Resistivity, Resistance, Electromotive Force and Circuits, Energy, and Power in Electric\par
Circuits, Theory of Metallic Conduction\par
\vspace{2pt}\textbf{- Direct-Current Circuits:}\par
Resistors in Series and Parallel, Kirchhoff’s Rules, Electrical Measuring Instruments, R-C\par
Circuits, Power Distribution Systems\par
\vspace{2pt}\textbf{- Magnetic Field and Magnetic Forces:}\par
Magnetism, Magnetic Field, Magnetic Field Lines, and Magnetic Flux, Motion of Charged\par
Particles in a Magnetic Field, Applications of Motion of Charged Particles, Magnetic Force on a\par
Current-Carrying Conductor, Force and Torque on a Current Loop, The Direct-Current Motor, The\par
Hall Effect\par
\vspace{2pt}\textbf{- Sources of Magnetic Field:}\par
Magnetic Field of a Moving Charge, Magnetic Field of a Current Element, Magnetic Field of a\par
Straight Current-Carrying Conductor, Force Between Parallel Conductors, Magnetic Field of a\par
Circular Current Loop, Ampere’s Law, Applications of Ampere’s Law, Magnetic Materials\par
\vspace{2pt}\textbf{- Electromagnetic Induction:}\par
Induction Experiments, Faraday’s Law, Lenz’s Law, Motional Electromotive Force, Induced\par
Electric Fields, Eddy Currents, Displacement Current and Maxwell’s Equations,\par
Superconductivity\par
\vspace{2pt}\textbf{- Inductance:}\par
Mutual Inductance, Self-Inductance, Inductors, Magnetic-Field Energy, The R-L Circuit, The LC Circuit, The L-R-C Series Circuit\par
\vspace{2pt}\textbf{- Alternating Current:}\par
Phasors and Alternating Currents, Resistance and Reactance, The L-R-C Series Circuit, Power in\par
Alternating-Current Circuits, Resonance in Alternating-Current Circuits, Transformers\par
\vspace{2pt}\textbf{- Electromagnetic Waves:}\par
Maxwell’s Equations and Electromagnetic Waves, Plane Electromagnetic Waves and the Speed\par
of Light, Sinusoidal Electromagnetic Waves, Energy and Momentum in Electromagnetic Waves,\par
Standing Electromagnetic Waves\par\vspace{4pt}
\textbf{TextBooks:}\par\vspace{2pt}
1. D. Holliday, R. Resnick and J. Walker, Fundamentals of Physics, 6th edition, John Wiley, 2000\par
\vspace{4pt}

\newpage
\textbf{Title of the Course:} Calculus 2\\
\textbf{Number of Credits:} 3\\
\textbf{Type of the Course:} Basic | Theoretical\\
\textbf{Pre-requisite:} Calculus 1\\
\textbf{Training hours:} 48h Theoretical\par\vspace{4pt}
\textbf{Course Objectives:}\par\vspace{2pt}
The main scope of this course is to teach the students some topics in introductory linear algebra
including matrix algebra and linear transformations and multivariable calculus including
multivariable functions, directional and partial derivatives, velocity and acceleration, tangent plane
and normal gradient line, cylindrical and spherical coordinates, vector field and line integrals,
surface integral, Green's theorem, Divergence theorem, and Stoke's theorem.\par\vspace{4pt}
\textbf{Syllabus of Courses:}\par\vspace{2pt}
Parametric equations; space coordinates; vector and space; numerical product; matrix 3x3 of threeindeterminate linear equations system; operation on lines; matrix reverse; solving equations\par
system; linear independence; base in R2; R3 linear transform and its matrix; determinant 3x3 and\par
characteristic value and vector; vector product; second order line and plane equations; two vector\par
functions and its derivative; speed and acceleration; bending; normal vector to a curve;\par
multivariable function; directional and partial derivative; tangent plane and normal line to a curve;\par
multivariable function; directional and partial derivative; tangent plane and projecting line of\par
gradient; chain of rule for partial derivative; exact differential; second kind and third kind integrals\par
and their application in geometrical and physical problems; transform of integration arrangement\par
(without accurate affirmation); cylindrical and spherical coordinates; vector field; curvilinear\par
integral; surface integral; divergence; curl; Laplacian; potential of green space and divergence and\par
stochastic.\par\vspace{4pt}
\textbf{TextBooks:}\par\vspace{2pt}
1. Tom M. Apostol, "Calculus Vol. 2: Multi-Variable Calculus and Linear Algebra with\par
Applications\par
to\par
Differential\par
Equations\par
and\par
Probability",\par
2012.\par
2. George B. Thomas and Ross L. Finney, "Calculus and Analytic Geometry (9th Edition)", 1995.\par
\vspace{4pt}

\newpage
\textbf{Title of the Course:} Discrete Mathematics\\
\textbf{Number of Credits:} 3\\
\textbf{Type of the Course:} Major | Theoretical\\
\textbf{Pre-requisite:} Calculus 1\\
\textbf{Training hours:} 48h Theoretical\par\vspace{4pt}
\textbf{Course Objectives:}\par\vspace{2pt}
Combinatorial reasoning underlies all analysis of a computer system and a discrete mathematics
course seeks to develop proficiency in basic discrete math problem-solving in the way that a
calculus course develops proficiency in basic analysis problem-solving. This course should teach
students to think logically and mathematically, reason and model combinatorially, and use discrete
structures. Some important themes interwoven in this course are mathematical reasoning and
mathematical induction, combinatorial analysis and enumeration, discrete structures, graphs, and
modeling the problems using them.\par\vspace{4pt}
\textbf{Syllabus of Courses:}\par\vspace{2pt}
\vspace{2pt}\textbf{- Induction and Recursion}\par
Mathematical Reasoning, Mathematical Induction, Strong Induction and Well-Ordering,\par
Recursive Definitions and Structural Induction, Recursive Algorithms\par
\vspace{2pt}\textbf{- Logics and Proofs}\par
Propositional Logic, Propositional Equivalences, Predicates and Quantifiers, Nested Quantifiers,\par
Rules of Inference, Introduction to Proofs, Proof Methods, and Strategy\par
\vspace{2pt}\textbf{- Counting}\par
The Basics of Counting, The Product Rule, The Sum Rule, the Tree Diagram, The Pigeonhole\par
Principle, Permutations and Combinations, Binomial Coefficients, Binomial Theorem, Pascal’s\par
Identities and Triange, Vander monde's Identity, Generalized Permutations and Combinations,\par
Permutations with Repetitions, Combinations with Repetitions, Distributing Objects into Boxes,\par
Algorithms For Generating Permutations and Combinations\par
\vspace{2pt}\textbf{- Advanced Counting Techniques}\par
Recurrence Relations, Modeling with Recurrence Relations, Fibonacci Numbers, The Tower of\par
Hanoi, Catalan Numbers, Solving Linear Recurrence Relations, Linear Homogeneous Recurrence\par
Relations with Constant Coefficients, Linear Nonhomogeneous Recurrence Relations with\par
Constant Coefficients, Divide-and-Conquer Algorithms, and Recurrence Relations, Formal Power\par
Series, Generating Functions, Counting Problems and Generating Functions, Inclusion-Exclusion,\par
Applications of Inclusion-Exclusion\par
\vspace{2pt}\textbf{- Relations}\par
Relations and Their Properties, n-ary Relations and Their Applications, Representing Relations,\par
Closures of Relations, Equivalence Relations, Partial Orderings.\par
\vspace{2pt}\textbf{- Graphs}\par
Graphs and Graph Models, Incidence and Adjacency Matrices of a Graph, Graph Terminology\par
and Special Types of Graphs, Complete graphs, Bipartite Graphs, Representing Graphs and Graph\par
Isomorphism, Connectivity, Euler and Hamilton Paths and Circuits, Algorithm to Construct Euler\par
Circuits, Ore’s Theorem, Dirac’s Theorem, Shortest-Path Problems, Dijkstra’s Algorithm, Planar\par
Graphs, Euler’s Formula, Kuratowski’s Theorem, Graph Coloring\par
\vspace{2pt}\textbf{- Trees}\par
Introduction to Trees, Rooted Trees, m-ary Trees, Applications of Trees, Huffman Coding, Tree\par
Traversal, Traversal Algorithms: Preorder, Inorder, Postorder, Spanning Trees, Depth First Search\par
and Breadth First Search Algorithms, Minimum Spanning Trees, Prim’s Algorithm, Kruskal’s\par
Algorithm\par\vspace{4pt}
\textbf{TextBooks:}\par\vspace{2pt}
1. R. P. Grimaldi, Discrete and Combinatorial Mathematics: An Applied Introduction, 5th Edition,\par
Pearson, 2017.\par
2. K. H. Rosen, Discrete Mathematics and Its Applications, 7th Edition, McGraw-Hill, 2011.\par
3. S. S. Epp, Discrete Mathematics with Applications, 4th Edition, Brooks Cole, 2010.\par
\vspace{4pt}

\newpage
\nopagebreak
\textbf{Title of the Course:} Islamic Thought 2 (Prophethood and Imamat)\\
\textbf{Number of Credits:} 2\\
\textbf{Type of the Course:} General | Theoretical\\
\textbf{Training hours:} 32h Theoretical\par\vspace{4pt}
\textbf{Course Objectives:}\par\vspace{2pt}
Becoming familiar with instructions in Islam, continuing "Islamic Thought I" discussions.
Expanding students' knowledge about religion, prophethood, Islam, Imamat, and leadership,
which is divided into three sections.\par\vspace{4pt}

\newpage
\vspace*{\fill}
\begin{center}
  {\fontsize{28pt}{34pt}\selectfont \textbf{SEMESTER 3}}\par\vspace{12pt}
  {\fontsize{13pt}{17pt}\selectfont Academic Year 2024–2025 (1st Semester)}
\end{center}
\vspace*{\fill}

\newpage
\textbf{Title of the Course:} Engineering Probability and Statistics\\
\textbf{Number of Credits:} 3\\
\textbf{Type of the Course:} Basic | Theoretical\\
\textbf{Pre-requisite:} Calculus 2\\
\textbf{Training hours:} 48h Theoretical\par\vspace{4pt}
\textbf{Course Objectives:}\par\vspace{2pt}
All branches of engineering need powerful analytic tools to deal with uncertainty, model the
data/information, and predict and analyze underlying systems. Many of the tools necessary for
these analyses have their foundations in probability and statistics. With this in mind, this course
provides elementary probabilistic and statistical concepts and the methods to apply them to
engineering problems.\par\vspace{4pt}
\textbf{Syllabus of Courses:}\par\vspace{2pt}
\vspace{2pt}\textbf{- Probability (Axioms, Interpretations, and Properties), Counting Techniques, Conditional}\par
Probability, Independence, Bayes’ Rule\par
\vspace{2pt}\textbf{- Discrete Random variables, Probability Mass Function, Probability Distribution Function,}\par
Expectation,\par
and\par
Variance\par
\vspace{2pt}\textbf{- Especial Discrete Distributions (Uniform, Binomial, Negative Binomial, HyperGeometric,}\par
Poisson)\par
\vspace{2pt}\textbf{- Continuous Random Variables, Probability Density Function, Probability Distribution Function,}\par
Expectation,\par
and\par
Variance,\par
Functions\par
of\par
a\par
Random\par
Variable\par
\vspace{2pt}\textbf{- Especial Continuous Distributions (Uniform, Exponential, Gamma, Beta, Weibull, Normal)}\par
\vspace{2pt}\textbf{- Joint Probability Distributions (Discrete and Continuous Cases), Conditional Distributions,}\par
Conditional Expectations, Independence\par
\vspace{2pt}\textbf{- Sampling Distributions, Distributions of the Sample Mean, and Sample Variance}\par
\vspace{2pt}\textbf{- Point Estimation: Maximum Likelihood Method, Method of Moments, Evaluating of Estimators}\par
(Biasedness, Variance, Mean of Squared Errors)\par
\vspace{2pt}\textbf{- Interval Estimation: Confidence Intervals for Mean, Variance, and proportion}\par
\vspace{2pt}\textbf{- Tests of Hypotheses based on a Single Sample (Tests for Mean, Variance, and Proportion, the}\par
Concept of P-Value)\par
\vspace{2pt}\textbf{- Tests of Hypotheses Based on Two Samples (Test for a Difference between Two Means, Test}\par
for a Difference between Two Proportions, Test for the ratio of Two Variances, T-Test for Paired\par
Data)\par
\vspace{2pt}\textbf{- Determination of Sample Size in Confidence Intervals and Testing Hypotheses}\par
\vspace{2pt}\textbf{- Nonparametric Tests (Goodness-of-Fit Test, Test of Independence in Two-Way Contingency}\par
Tables)\par
\vspace{2pt}\textbf{- Regression (Simple Linear Regression, Normal Regression Analysis, Correlation)}\par\vspace{4pt}
\textbf{TextBooks:}\par\vspace{2pt}
1. Walpole and Mayers, Probability and Statistics for Engineering and Scientists, 6th\par
edition, Prentice Hall. 1998\par
2. R.V. Hogg and T. Elliot, Probability and Statistical Inference, 4th edition, Mac Millan, 1993\par
3. Guttman, I., Wilks, S.S., Hunter, J.S. Introductory Engineering Statistics, J. Wiley, 3rd Ed.,\par
1982.\par
\vspace{4pt}

\newpage
\textbf{Title of the Course:} Logic Circuits\\
\textbf{Number of Credits:} 3\\
\textbf{Type of the Course:} Specialized | Theoretical\\
\textbf{Training hours:} 48h Theoretical\par\vspace{4pt}
\textbf{Course Objectives:}\par\vspace{2pt}
This course covers the design and analysis of digital logic circuits including combinatorial and
sequential logic circuits.\par\vspace{4pt}
\textbf{Syllabus of Courses:}\par\vspace{2pt}
Number systems, binary systems, binary codes, addition, subtraction, multiplication and division,\par
Boolean algebra, theorems, logic functions, sum-of-products form, product-of-sums form,\par
Karnaugh maps, logic minimization using K-maps, logic gates (RTL, DTL, TTL) and related\par
topics (FAN-IN, FAN-OUT), combinatorial circuits design, multiplexers, decoders, comparators,\par
sequential circuits design, flip-flops, shift registers, counters, synchronous and asynchronous logic\par
circuits, finite state machines, the timing of sequential circuits\par
Programming language: Verilog\par\vspace{4pt}
\textbf{TextBooks:}\par\vspace{2pt}
1. Morris Mano, Michael Ciletti, Digital Design. 4th ed. Prentice Hall, 2006.\par
2. Victor Nelson, Troy Nagle, Bill Carroll, David Irwin, Digital Logic Circuit Analysis and Design.\par
1st ed. Prentice Hall, 1995.\par
\vspace{4pt}

\newpage
\textbf{Title of the Course:} Electric Circuits\\
\textbf{Number of Credits:} 3\\
\textbf{Type of the Course:} Specialized | Theoretical\\
\textbf{Training hours:} 48h Theoretical\par\vspace{4pt}
\textbf{Syllabus of Courses:}\par\vspace{2pt}
Coulomb's, Ohm's, and Kirchhoff’s laws; electrostatics and electromagnetics; resistance,\par
capacitance, inductance; series and parallel circuits, independent and dependent voltage and\par
current sources; energy, power; Superposition, Thevenin, and Norton Theorems; DC and AC;\par
initial, steady state and transient conditions; frequency selective circuits and resonance; polyphase circuits and transformers.\par\vspace{4pt}
\textbf{TextBooks:}\par\vspace{2pt}
1- W. H. Hayt, J. E. Kemmerly, and S. M. Durbin, Engineering Circuit Analysis, 8th Edition,\par
McGraw Hill, 2011.\par
2- C. K. Alexander, M. N. O. Sadiku, Fundamentals of Electrical Circuits, 5th Edition, McGraw\par
Hill, 2012\par
3- A. Agarwal, J. H. Lang, Foundations of Analog and Digital Electronic Circuits, Morgan\par
Kaufmann Publishers, Elsevier, 2005\par
\vspace{4pt}

\newpage
\nopagebreak
\textbf{Title of the Course:} Islamic Revolution in Iran\\
\textbf{Number of Credits:} 2\\
\textbf{Type of the Course:} General | Theoretical\\
\textbf{Training hours:} 32h Theoretical\par\vspace{4pt}
\textbf{Course Objectives:}\par\vspace{2pt}
Theoretical acquaintance with the causes and factors of the Islamic Revolution and analytical study
of cultural, social, and political developments of the Islamic Revolution and subsequent issues.\par\vspace{4pt}

\par\vspace{14pt}
\nopagebreak
\textbf{Title of the Course:} Family Schematization and Population\\
\textbf{Number of Credits:} 2\\
\textbf{Type of the Course:} General | Theoretical\\
\textbf{Training hours:} 32h Theoretical\par\vspace{4pt}
\textbf{Course Objectives:}\par\vspace{2pt}
Familiarity of students with the basics of family knowledge and how to form, consolidate, and
excel it as demographic characteristics and their role in improving the quality of life.\par\vspace{4pt}

\newpage
\textbf{Title of the Course:} Data Structures\\
\textbf{Number of Credits:} 3\\
\textbf{Type of the Course:} Major | Theoretical\\
\textbf{Pre-requisite:} Advanced Programming\\
\textbf{Training hours:} 48h Theoretical\par\vspace{4pt}
\textbf{Course Objectives:}\par\vspace{2pt}
Familiarity with Information Structures, Effects of Structures in Produced Applications, Selecting
Optimized Inside-Memory Structures, and Organizing Memory based on requirements.\par\vspace{4pt}
\textbf{Syllabus of Courses:}\par\vspace{2pt}
Arrays, Vectors, matrices and their application, sparse matrices, heap and stack, queue, lists, linked\par
lists (linear, cyclic, doubly and multiply linked lists), introduction to trees, binary trees (decisionmaking trees, gaming trees, search trees and …), graphs (representation, sweeping and\par
applications), Spanning trees, Dynamic memory allocation methods, searching, sorting and\par
merging algorithms.\par\vspace{4pt}
\textbf{TextBooks:}\par\vspace{2pt}
1. Thomas H. Cormen, Charles E. Leiserson, Ronald L. Rivest, Clifford Stein, Introduction to\par
Algorithms, Third Edition, The MIT Press, Cambridge, Massachusetts 2009\par
2. A.V. Aho, J.E. Hopcroft, J.D. Ullman, Data Structures and Algorithms, Addison Wesley, 1983\par
\vspace{4pt}

\newpage
\vspace*{\fill}
\begin{center}
  {\fontsize{28pt}{34pt}\selectfont \textbf{SEMESTER 4}}\par\vspace{12pt}
  {\fontsize{13pt}{17pt}\selectfont Academic Year 2024–2025 (2nd Semester)}
\end{center}
\vspace*{\fill}

\newpage
\textbf{Title of the Course:} Logic Circuits Laboratory\\
\textbf{Number of Credits:} 1\\
\textbf{Type of the Course:} Specialized | Practical\\
\textbf{Pre-requisite:} Logic Circuits\\
\textbf{Training hours:} 32h Practical\par\vspace{4pt}
\textbf{Course Objectives:}\par\vspace{2pt}
The logic circuits laboratory is provided with the logic circuits course and its purpose is to acquaint
students with the practical components of logic circuits and how these circuits work and to acquire
skills in the development of circuits and digital systems.\par\vspace{4pt}
\textbf{TextBooks:}\par\vspace{2pt}
1- C. Roth, L. Kinney, Fundamentals of Logic Design, 7th Edition, Cengage Learning, 2014.\par
2- F. Vahid, Digital Design with RTL Design, VHDL, and Verilog, 2nd Edition, Wiley, 2011.\par
\vspace{4pt}

\newpage
\textbf{Title of the Course:} Engineering Mathematics\\
\textbf{Number of Credits:} 3\\
\textbf{Type of the Course:} Specialized | Theoretical\\
\textbf{Pre-requisite:} Calculus 2\\
\textbf{Training hours:} 48h Theoretical\par\vspace{4pt}
\textbf{Course Objectives:}\par\vspace{2pt}
This course provides mathematical techniques in the more advanced areas of mathematics that are
of most relevance to engineering disciplines.\par\vspace{4pt}
\textbf{Syllabus of Courses:}\par\vspace{2pt}
\vspace{2pt}\textbf{- Fourier Series, Integrals, and The Fourier Transform}\par
Introduction to Fourier Series, Convergence of Fourier Series and Their Integration and\par
Differentiation, Fourier Sine and Cosine Series, Other Forms of Fourier Series, Frequency and\par
Amplitude Spectra of a Function, Double Fourier Series. The Fourier Integral, The Fourier\par
Transform, Fourier Cosine and Sine Transforms, Sturm–Liouville Problems, Eigenfunctions, and\par
Orthogonality\par
\vspace{2pt}\textbf{- Partial Differential Equations}\par
The Method of Characteristics, Wave Propagation and First Order PDEs, Generalizing Solutions,\par
Conservation Laws and Shocks, The Three Fundamental Types of Linear Second Order PDE,\par
Classification and Reduction to Standard Form of a Second Order Constant Coefficient Partial\par
Differential Equation, Boundary Conditions and Initial Conditions, Waves and the OneDimensional Wave Equation, The D’Alembert Solution of the Wave Equation and Applications,\par
Separation of Variables, Some General Results for the Heat and Laplace Equation, An\par
Introduction to Laplace and Fourier Transform Methods for PDEs\par
\vspace{2pt}\textbf{- Complex Analysis}\par
Complex Functions and Mappings, Limits, Derivatives, and Analytic Functions, Harmonic\par
Functions and Laplace’s Equation, Elementary Functions, Inverse Functions and Branches,\par
Complex Integrals, Contours, the Cauchy–Goursat Theorem, and Contour Integrals, The Cauchy\par
Integral Formulas, Some Properties of Analytic Functions, Complex Power Series and Taylor\par
Series, Uniform Convergence, Laurent Series and the Classification of Singularities, Residues and\par
the Residue Theorem, Evaluation of Real Integrals using Residues, The Inversion Integral for the\par
Laplace Transform, Conformal Mapping, Conformal Mapping, and Boundary Value Problems\par
\vspace{2pt}\textbf{- Calculus of Variations}\par
Maximum or minimum of an Integral, Euler’s equations, Lagrange’s equations, Sturm–Liouville\par
systems, and Rayleigh-Ritz methods.\par\vspace{4pt}
\textbf{TextBooks:}\par\vspace{2pt}
1. Michael Greenberg, Advanced Engineering Mathematics, 2nd edition, 1998\par
2. Erwin Kreyszig, Advanced Engineering Mathematics, tenth edition, 2011\par
\vspace{4pt}

\newpage
\textbf{Title of the Course:} Computer Workshop\\
\textbf{Number of Credits:} 3\\
\textbf{Type of the Course:} Basic | Practical\\
\textbf{Pre-requisite:} Introduction to Computing Systems and Programming\\
\textbf{Training hours:} 48h Practical\par\vspace{4pt}
\textbf{Course Objectives:}\par\vspace{2pt}
Familiarity with accessory systems such as card readers, printers, magnetic tape, discs, and
consoles, manner of work with terminals, familiarity with compilers and editors, familiarity with
JCL language, familiarity with computer organization of a center, familiarity with prepared
software packages such as database, spreadsheet, lotus, familiarity with the important programs of
a system such as a sort, merge, creation and copy of files, etc., familiarity with usage manner of
an operating system of microcomputer\par\vspace{4pt}

\newpage
\textbf{Title of the Course:} Artificial Intelligence\\
\textbf{Number of Credits:} 3\\
\textbf{Type of the Course:} Specialized | Theoretical\\
\textbf{Training hours:} 48h Theoretical\par\vspace{4pt}
\textbf{Course Objectives:}\par\vspace{2pt}
AI is an introductory course in Artificial Intelligence. The goal is to acquire knowledge of
intelligent systems and agents, formalization of knowledge, reasoning with and without
uncertainty, machine learning, and applications at a basic level.\par\vspace{4pt}
\textbf{Syllabus of Courses:}\par\vspace{2pt}
Intelligent agents | Problem-solving, complexity metrics | Uninformed search (BFS, DFS, Uniform\par
cost, depth-limited, iterative deepening) | Informed search (Greedy, A*) | Another search\par
(simulated annealing, genetic algorithms) | Game playing (minimax, alpha-beta pruning,\par
nondeterministic games) | Knowledge-based agents | Propositional logic, syntax, semantics,\par
inference, limitations | First-order logic (including inference, unification, Goedel’s theorem,\par
generalized modus ponens, resolution, forward chaining, backward chaining, situation calculus) |\par
Building a knowledge base (ontologies, translation between English and FOL, difficult cases like\par
events, substances, composite objects) | Logical reasoning systems (Implementing indexing,\par
retrieval, and unification. Theorem provers. Frame systems and semantic networks) | Planning\par
(including STRIPS) | Uncertainty | Decision Trees | Fuzzy logic | Neural networks (including\par
Hopfield, perceptrons, backpropagation)\par\vspace{4pt}
\textbf{TextBooks:}\par\vspace{2pt}
1. Artificial Intelligence: a modern approach, S. Russell and P. Norvig, Prentice Hall.\par
\vspace{4pt}

\newpage
\textbf{Title of the Course:} Computer Architecture\\
\textbf{Number of Credits:} 3\\
\textbf{Type of the Course:} Specialized | Theoretical\\
\textbf{Pre-requisite:} Logic Circuits\\
\textbf{Training hours:} 48h Theoretical\par\vspace{4pt}
\textbf{Course Objectives:}\par\vspace{2pt}
This course aims to give an understanding of the mechanisms for implementing the programmer's
idealized computer. It builds on the introduction to hardware and simple processors in the Digital
Systems course. The Computer Architecture course aims to describe a broad range of architectural
designs and to contrast them, highlighting the design decisions they incorporate, and how these
design decisions impact program performance.\par\vspace{4pt}
\textbf{Syllabus of Courses:}\par\vspace{2pt}
C programming language; x86-64 assembly language; Machine representation of programs;\par
Control structures for processors, register transfer level description of hardware; Hardware\par
description languages and simulation in Verilog; Instruction set design, instruction formats,\par
addressing modes, ISAs; A sequential x86-64 design; Processor pipelining, pipeline hazard\par
detection, stalling and forwarding; Out-of-order execution and program performance optimization;\par
RISC and CISC instruction sets (x86-64 and MIPS); Memory hierarchy and memory caches;\par
Microcontrollers for embedded applications; Alternative architectures.\par\vspace{4pt}
\textbf{TextBooks:}\par\vspace{2pt}
1. D A Patterson \& J L Hennessy, Computer Organization and Design: The hardware/software\par
interface, Morgan-Kaufmann (Fifth edition) 2013.\par
\vspace{4pt}

\newpage
\textbf{Title of the Course:} Machines and Language Theory\\
\textbf{Number of Credits:} 3\\
\textbf{Type of the Course:} Specialized | Theoretical\\
\textbf{Training hours:} 48h Theoretical\par\vspace{4pt}
\textbf{Course Objectives:}\par\vspace{2pt}
A presentation of the theory of computing, including coverage of the theory of formal languages
and automata, compatibility, computational complexity, and deterministic parsing of context-free
languages.\par\vspace{4pt}
\textbf{TextBooks:}\par\vspace{2pt}
1. Sipser. Introduction to the Theory of Computation. PWS Publishing company, 1997.\par
2. Linz, An Introduction to Formal Languages and Automata. 3rd edition ed., Jone,s and Bartlett\par
Publishers, 2001.\par
\vspace{4pt}

\newpage
\textbf{Title of the Course:} Electrical Measurement and Circuit Laboratory\\
\textbf{Number of Credits:} 1\\
\textbf{Type of the Course:} Specialized | Practical\\
\textbf{Pre-requisite:} Electric Circuits\\
\textbf{Training hours:} 32h Practical\par\vspace{4pt}
\textbf{Course Objectives:}\par\vspace{2pt}
In this laboratory, students are familiar with the electronic elements of resistors, capacitors, and
inductors, and the operation of electrical measuring devices such as multimeters and oscilloscopes,
as well as other devices used in electrical and electronic circuits such as power supply and signal
generators, as well as measuring parameters such as current. They learn voltage, frequency,
periodicity, phase difference, and power.\par\vspace{4pt}

\newpage
\textbf{Title of the Course:} Foundations of Data Science\\
\textbf{Number of Credits:} 3\\
\textbf{Type of the Course:} Elective | Theoretical\\
\textbf{Training hours:} 48h Theoretical\par\vspace{4pt}
\textbf{Course Objectives:}\par\vspace{2pt}
Introduction to the data science lifecycle and the principles of collecting, processing, analyzing, visualizing, and modeling data. Developing an understanding of statistical inference, exploratory data analysis, predictive and unsupervised learning, model evaluation, and responsible data-driven decision-making. Familiarity with modern data systems, machine learning operations, neural networks, large language models, and interpretable AI.\par\vspace{4pt}
\textbf{Syllabus of Courses:}\par\vspace{2pt}
Data science lifecycle; sampling, sampling bias, observational and experimental studies, confounding variables; statistical charts; probability theory, conditional probability, Bayes theorem, random variables, expectation and variance, normal and binomial distributions, Monte Carlo simulation; central limit theorem, confidence intervals, hypothesis testing and statistical power; data visualization design principles, preattentive attributes, dashboards and storytelling; data preprocessing, exploratory data analysis and feature engineering; SQL, NoSQL, graph databases and MongoDB; big data systems including Hadoop, Spark, Kafka, HBase and Cassandra; modeling principles, prediction versus explanation, training and evaluation splits; linear regression, loss functions and gradient descent; logistic regression; decision trees, random forests, boosting, support vector machines and k-nearest neighbors; regression, classification and clustering evaluation; MLOps, Docker, Kubernetes, MLflow, data and model drift, retraining; unsupervised learning, dimensionality reduction and k-means clustering; feedforward neural networks, convolutional neural networks and recurrent neural networks; image classification and time series; natural language processing, embeddings, transformers and large language models; retrieval-augmented generation and AI agents; model interpretability and explainability; privacy and ethics in data science.\par\vspace{4pt}
\textbf{TextBooks:}\par\vspace{2pt}
1. Gareth James, Daniela Witten, Trevor Hastie, and Robert Tibshirani, An Introduction to Statistical Learning: with Applications in R/Python (ISLR).\par
2. Trevor Hastie, Robert Tibshirani, and Jerome Friedman, The Elements of Statistical Learning (ESL).\par
3. Jake VanderPlas, Python Data Science Handbook: Essential Tools for Working with Data.\par
4. Joel Grus, Data Science from Scratch: First Principles with Python.\par
5. Aurélien Géron, Hands-On Machine Learning with Scikit-Learn, Keras, and TensorFlow.\par
\vspace{4pt}

\newpage
\vspace*{\fill}
\begin{center}
  {\fontsize{28pt}{34pt}\selectfont \textbf{SEMESTER 5}}\par\vspace{12pt}
  {\fontsize{13pt}{17pt}\selectfont Academic Year 2025–2026 (1st Semester)}
\end{center}
\vspace*{\fill}

\newpage
\textbf{Title of the Course:} Operating Systems Laboratory\\
\textbf{Number of Credits:} 1\\
\textbf{Type of the Course:} Specialized | Practical\\
\textbf{Training hours:} 32h Practical\par\vspace{4pt}
\textbf{Course Objectives:}\par\vspace{2pt}
The goal of this course is to create the right understanding and to learn the techniques of designing
and implementing system programs and operating systems based on an open-source operating
system. Students will get familiar with the implementation of basic design concepts of the
operating system, such as file management, process management, kernel-level programming, and
process synchronization in an open-source operating system, and learn to install and manage an
open-source operating system.
Upon successful completion of the course, students will be able to
1. Compile, install, and manage an open-source operating system
2. Implement the basic concepts of designing operating systems in an open-source operating
system
3. Program in kernel-level\par\vspace{4pt}
\textbf{Syllabus of Courses:}\par\vspace{2pt}
1) Introduction: An overview of Linux including its history, versions, Posix, Graphical User\par
Interface\par
2) Install and compile a version of Linux: Disk partition, compile, boot load, and use the\par
application menu and system.\par
3) Linux file system\par
4) Standard and advanced Shell\par
5) Managing Processes and Threads\par
6) Processor scheduling in Linux\par
7) Deadlock synchronization\par
8) Memory management\par\vspace{4pt}

\newpage
\textbf{Title of the Course:} Operating System\\
\textbf{Number of Credits:} 3\\
\textbf{Type of the Course:} Specialized | Theoretical\\
\textbf{Pre-requisite:} Computer Architecture\\
\textbf{Training hours:} 48h Theoretical\par\vspace{4pt}
\textbf{Course Objectives:}\par\vspace{2pt}
The course provides computer engineering students to understand the ways that applications
interact with the computer hardware as well as different methods of resource management and
their respective algorithms. Further, they will learn of basic mechanisms which are used in the
low-level management of computer systems. Such a student will also learn how to design a system
with the consideration of both hardware and software characteristics and limitations and improve
the quality of programs. The course is a combination of both practice and theory.
Upon successful completion of the course, students will have good insights
1. in different types of computer systems and applications,
2. to design, build, and manage system software,
3. to better discover the reasons for the degraded performance of computer systems and solve the
problem,
4. to propose and implement resource management policies.\par\vspace{4pt}
\textbf{Syllabus of Courses:}\par\vspace{2pt}
1) Introduction (Basic terminologies and concepts as well as different computer systems\par
architectures and software systems, Interrupts, Dual-mode, …)\par
2) Operating system structures (Different architectures for operating systems, system calls, APIs,\par
virtualization, and cloud computing)\par
3) Processes (PCB, context switch, short-term, medium-term, and long-term schedulers, process\par
creation and communication)\par
4) Threads (Multithreading models, thread management)\par
5) Process synchronization (Critical sections and respective solutions, semaphores, monitors,\par
classic synchronization problems)\par
6) CPU scheduling (Preemptive and non-preemptive scheduling, CPU scheduling algorithms,\par
SMP, SMT, multicore scheduling, real-time scheduling)\par
7) Deadlock handling (Necessary conditions for deadlock, resource-allocation graph, deadlock\par
prevention, avoidance, detection, and recovery)\par
8) Main memory management (Binding, physical and logical addresses, fragmentation, paging,\par
TLB, shared pages, segmentation)\par
9) Virtual memory management (Demand paging, page fault handling, copy-on-write, page\par
replacement policies, Belady’s anomaly, frame allocation, thrashing, memory-mapped I/O and\par
files, kernel memory management)\par
10) Storage management (File system structure, virtual file system, storage allocation methods)\par
11) I/O management (Interrupts, I/O subsystem, I/O management and performance aspects)\par\vspace{4pt}
\textbf{TextBooks:}\par\vspace{2pt}
1. W. Stallings, Operating Systems, 4th Edition, Prentice Hall, 2001\par
2. Silberschatz, J.L. Peterson, Operating Systems Concepts, Addison Wesley, 2000 Course\par
\vspace{4pt}

\newpage
\textbf{Title of the Course:} Computer Aided Digital System\\
\textbf{Number of Credits:} 3\\
\textbf{Type of the Course:} Specialized | Theoretical\\
\textbf{Training hours:} 48h Theoretical\par\vspace{4pt}
\textbf{Course Objectives:}\par\vspace{2pt}
Students in the Computer Aided Digital System course will implement and test digital circuits
ranging in complexity from basic logic gates to state machines that perform useful functions like
calculations, counting, timing, and a host of other applications.\par\vspace{4pt}

\newpage
\textbf{Title of the Course:} Signals and Systems\\
\textbf{Number of Credits:} 3\\
\textbf{Type of the Course:} Specialized | Theoretical\\
\textbf{Pre-requisite:} Engineering Mathematics\\
\textbf{Training hours:} 48h Theoretical\par\vspace{4pt}
\textbf{Course Objectives:}\par\vspace{2pt}
Introduction to the fundamentals of continuous and discrete time signal and system analysis. This
course will cover linear system analysis including impulse response and convolution, Fourier
series, Fourier transforms, sampling, discrete-time signal and system analysis, and Z-transforms.
The student is expected to demonstrate the ability to:
- Apply the convolution theorem for continuous-time signals
- Evaluate the Fourier Series of periodic signals
- Determine the Fourier Transform of energy signals
- Make use of Fourier Transform Properties
- Analyze a discrete-time LTI system using discrete linear convolution
- Use Z-transform for analyzing discrete-time signals and systems
- Convert a continuous-time signal to the discrete-time domain and reconstruct using the
sampling theorem\par\vspace{4pt}
\textbf{Syllabus of Courses:}\par\vspace{2pt}
\vspace{2pt}\textbf{- Introduction}\par
Motivation: Why signal processing is important, Mathematical preliminaries\par
\vspace{2pt}\textbf{- Signal Basics}\par
Types of signals, Elementary signals\par
\vspace{2pt}\textbf{- Time-Domain Analysis of Continuous-Time Systems}\par
Types of systems, Linear time-invariant (LTI) systems, Convolution\par
\vspace{2pt}\textbf{- Continuous-Time Fourier Series (CTFS)}\par
CTFS derivation, CTFS properties, Convergence\par
\vspace{2pt}\textbf{- Continuous-Time Fourier Transform (CTFT)}\par
CTFT derivation, CTFT properties, CTFT, convolution, and LTI systems, Convergence\par
\vspace{2pt}\textbf{- Laplace Transform}\par
Region of convergence, Implementing continuous-time systems\par
\vspace{2pt}\textbf{- Time-Domain Analysis of Discrete-Time Systems}\par
Types of systems, Linear time-invariant (LTI) systems, Convolution\par
\vspace{2pt}\textbf{- Time-Domain Analysis of Discrete-Time Systems}\par
Types of systems, Linear time-invariant (LTI) systems, Convolution\par
\vspace{2pt}\textbf{- Discrete-Time Fourier Transform (DTFT)}\par
DTFT derivation, DTFT properties, DTFT, convolution, and LTI systems, Convergence\par
\vspace{2pt}\textbf{- Discrete Fourier Transform (DFT)}\par
DFT, convolution, and LTI systems, Convergence, Fast Fourier Transform (FFT)\par
\vspace{2pt}\textbf{- Sampling and Reconstruction}\par
Sampling and the Nyquist theorem, Sinc interpolation, Practical reconstruction\par
\vspace{2pt}\textbf{- Z-Transform}\par
Region of convergence, Implementing discrete-time systems\par\vspace{4pt}
\textbf{TextBooks:}\par\vspace{2pt}
1. Alan V. Oppenheim, A. S. Willsky and S. Hamid Nawab, Signals and Systems, 2nd ed.,\par
Prentice-Hall, 1996.\par
2. S. Haykin and B. Van Veen, Signals and Systems, 2nd ed., Wiley, 2003.\par
\vspace{4pt}

\newpage
\nopagebreak
\textbf{Title of the Course:} English Language\\
\textbf{Number of Credits:} 3\\
\textbf{Type of the Course:} General | Theoretical\\
\textbf{Training hours:} 48h Theoretical\par\vspace{4pt}
\textbf{Course Objectives:}\par\vspace{2pt}
General Language courses, offered at all levels of the Intensive Language Program, seek to
improve students' Language skills in a great variety of discourse settings. Language skills
addressed include listening, fluency development, oral intelligibility, reading, grammar, writing,
and vocabulary development.\par\vspace{4pt}

\newpage
\textbf{Title of the Course:} Algorithm Design 1\\
\textbf{Number of Credits:} 3\\
\textbf{Type of the Course:} Major | Theoretical\\
\textbf{Pre-requisite:} Data Structures\\
\textbf{Training hours:} 48h Theoretical\par\vspace{4pt}
\textbf{Course Objectives:}\par\vspace{2pt}
The main purpose of this course is to teach the basic concepts and common methods of analysis
and design of algorithms. In this lesson, students learn how to analyze a given problem and find
some possible algorithms to solve it. Then analyze and compare those algorithms in terms of
computational complexity and identify the best ones based on the characteristics of the input
samples. In this lesson, basic algorithms for solving some common problems will be introduced.\par\vspace{4pt}
\textbf{Syllabus of Courses:}\par\vspace{2pt}
Review of essential points of Data Structures, problem-solving methods (for each method some\par
problems and a special algorithm for the problem shall be represented and analyzed), introduction\par
to complexity, divide \& conquer method (problems to be analyzed: max and min of an array,\par
multiplication of two n-bit number, Strassen method about matrix multiplication, round robin\par
algorithm, sorting with Quicksort algorithm), dynamic programming method (matrix\par
multiplication, traveling salesman problem, polygon triangulation), greedy algorithms(scheduling\par
problem, Huffman Code, making change), methods based on exhaustive search, alpha-beta\par
pruning (puzzle, tic-tac-tac), the revelation of methods for problems, graph algorithms(searching\par
methods of the graph, Dijkstra algorithms, minimum spanning tree, Floyd algorithm, topological\par
sorting and…), maximum flow networks and other problems.\par\vspace{4pt}
\textbf{TextBooks:}\par\vspace{2pt}
1. Thomas H. Cormen, Charles E. Leiserson, Ronald L. Rivest, Clifford Stein, Introduction to\par
Algorithms, Third Edition, The MIT Press, Cambridge, Massachusetts 2009.\par
2. R.E. Neapolitan and K.Naimipour, Foundations of Algorithms using C++ Pseudo Codes, 2nd\par
Edition,\par
Johns\par
and\par
Bartlett\par
Publishers,\par
1998.\par
\vspace{4pt}

\newpage
\textbf{Title of the Course:} Basics of Machine Learning\\
\textbf{Number of Credits:} 3\\
\textbf{Type of the Course:} Elective | Theoretical\\
\textbf{Training hours:} 48h Theoretical\par\vspace{4pt}
\textbf{Course Objectives:}\par\vspace{2pt}
Understanding fundamental concepts of machine learning and pattern recognition, including probabilistic modeling, statistical decision theory, parameter estimation, supervised and unsupervised learning. Developing the ability to formulate classification and regression problems, select and optimize learning models, and apply neural networks, feature selection, dimensionality reduction, and clustering methods.\par\vspace{4pt}
\textbf{Syllabus of Courses:}\par\vspace{2pt}
Introduction to machine learning concepts and fundamentals of data; Bayesian decision theory, Bayesian classifiers, risk minimization and decision boundaries; Naive Bayes classification; maximum likelihood estimation (MLE) and maximum a posteriori (MAP) estimation; optimization methods; linear, logistic and nonlinear regression; decision trees; support vector machines (SVM); introduction to neural networks and their applications; introduction to deep vision; principal component analysis (PCA), linear discriminant analysis (LDA), and feature selection; clustering and unsupervised learning.\par\vspace{4pt}
\textbf{TextBooks:}\par\vspace{2pt}
1. Christopher M. Bishop, Pattern Recognition and Machine Learning, Springer.\par
2. Richard O. Duda, Peter E. Hart, and David G. Stork, Pattern Classification, Wiley, 2001.\par
3. Sergios Theodoridis and Konstantinos Koutroumbas, Pattern Recognition, Academic Press, 2009.\par
4. Trevor Hastie, Robert Tibshirani, and Jerome Friedman, The Elements of Statistical Learning: Data Mining, Inference, and Prediction, Second Edition.\par
\vspace{4pt}

\newpage
\vspace*{\fill}
\begin{center}
  {\fontsize{28pt}{34pt}\selectfont \textbf{SEMESTER 6}}\par\vspace{12pt}
  {\fontsize{13pt}{17pt}\selectfont Academic Year 2025–2026 (2nd Semester)}
\end{center}
\vspace*{\fill}

\newpage
\textbf{Title of the Course:} Systems Analysis and Design\\
\textbf{Number of Credits:} 3\\
\textbf{Type of the Course:} Elective | Theoretical\\
\textbf{Training hours:} 48h Theoretical\par\vspace{4pt}
\textbf{Course Objectives:}\par\vspace{2pt}
Just as solving a mathematical problem requires first identifying its complexities and dimensions,
designing appropriate strategies to solve the problem, and then proceeding to solve it, such as
building computer systems for a set of people (which in many cases may be more complex
mathematical problems) also need to be analyzed and designed before implementation. In this
course, students are introduced to analysis and design skills such as requirements analysis,
feasibility analysis, modeling, process analysis, architectural design, and UI / UX. In addition to
these skills, this course will include computer systems management and production planning skills
such as production methods, process automation, and project management concepts in such
projects.\par\vspace{4pt}
\textbf{TextBooks:}\par\vspace{2pt}
1. L.D. Bentley, K.C. Dittman, and J.L. Whitten. Systems analysis and design methods.\par
Irwin/McGraw Hill, 2007.\par
2. K.S. Rubin. Essential Scrum: A practical guide to the most popular Agile process. AddisonWesley, 2012.\par
3. M. Fowler, C. Kobryn, and K. Scott. UML distilled: A brief guide to the standard object\par
modeling language. Addison-Wesley Professional, 2004.\par
\vspace{4pt}

\newpage
\textbf{Title of the Course:} Technical English\\
\textbf{Number of Credits:} 2\\
\textbf{Type of the Course:} Specialized | Theoretical\\
\textbf{Training hours:} 32h Theoretical\par\vspace{4pt}
\textbf{Course Objectives:}\par\vspace{2pt}
This subject aims at raising students’ specific language ability in reading and writing academic
texts of their major disciplines. The subject will use reading texts from chapters of books or journal
articles recommended by teachers of different majors for reading comprehension. These texts will
also be used for analysis to enable students to develop an awareness of the genre in that particular
discipline.\par\vspace{4pt}

\newpage
\textbf{Title of the Course:} Computer Networks\\
\textbf{Number of Credits:} 3\\
\textbf{Type of the Course:} Specialized | Theoretical\\
\textbf{Training hours:} 48h Theoretical\par\vspace{4pt}
\textbf{Course Objectives:}\par\vspace{2pt}
The purpose of this course is to become familiar with the principles and practice of data
networking; learn how to write applications that use the network; get some understanding of
network internals in a hands-on way; explain the OSI model; recognize the primary network
protocols and distinguish between centralized, client/server, and peer-to-peer systems; and
distinguish between Local Area Networks (LANs) and Wide Area Networks (WANs) and identify
the components used to expand a LAN into a WAN.
Upon successful completion of the course, students will be able to:
1. To master the terminology and concepts of the OSI reference model and the TCP‐IP reference
model.\par\vspace{4pt}
\textbf{Syllabus of Courses:}\par\vspace{2pt}
Review of fundamental concepts in networking and the Internet\par
1. Internet protocols and components in the edge and the core of the network\par
2. Performance concepts such as queue, delay, loss, and throughput\par
3. Layered architecture\par
Application Layer\par
1. Web and HTTP protocol\par
2. Electronic Mail and SMTP protocol\par
3. Internet’s Directory Service and DNS protocol\par
4. P2P applications and P2P file distribution\par
5. Video streaming and DASH protocol\par
6. Content distribution and Cache\par
7. Socket programming\par
Transport Layer\par
1. Multiplexing and demultiplexing\par
2. UDP protocol\par
3. Principles of reliable transfer and GBN and SR protocols\par
4. TCP protocol, round-trip time estimation, and flow control\par
5. Principles of congestion control\par
Network Layer: Data Plane\par
1. Data and control planes\par
2. Router architecture\par
3. Addressing, IPv4 and IPv6\par
4. Network address translation and NAT protocol\par
5. Generalized forwarding (SDN) and OpenFlow protocol\par
Network Layer: Control Plane\par
1. LS and DV routing algorithms\par
2. Intra-AS routing and OSPF protocol\par
3. Inter-AS routing and BGP protocol\par
4. SDN control plane\par
5. ICMP and SNMP protocols\par
Data Link Layer and LANs\par
1. The data link layer services\par
2. Error-Detection and –Correction techniques\par
3. Multiple access protocols\par
4. Local area networks, Ethernet, ARP protocol, and VLAN protocol\par
5. Link virtualization and MPLS protocol\par
6. Data center networking\par\vspace{4pt}
\textbf{TextBooks:}\par\vspace{2pt}
1. Computer Networking: a top-down approach / James F. Kurose, Keith W. Ross.6th ed, 2013. 2.\par
Computer networks / Andrew S. Tanenbaum, David J. Wetherall. 5th ed, 2011.\par
\vspace{4pt}

\newpage
\textbf{Title of the Course:} Database Design\\
\textbf{Number of Credits:} 3\\
\textbf{Type of the Course:} Elective | Theoretical\\
\textbf{Training hours:} 48h Theoretical\par\vspace{4pt}
\textbf{Course Objectives:}\par\vspace{2pt}
This course introduces relational database systems. In this course, the students will learn how to
design and create relational databases and how to use them. They will also learn about the internals
of database management systems: how such systems store data, optimize and execute queries, and
process transactions.
Upon successful completion of the course, students will be able
1. to understand the basics of a database system and its architecture
2. to design a database for an application: draw an entity-relationship (ER) diagram from a problem
specification, convert the ER diagram to relations, and normalize the relations
3. to write database queries in SQL
4. to write relational algebra and relational calculus expressions for queries
5. to understand how the relations and indexes are stored in a database system
6. to identify useful indexes for a specific database and workload
7. to understand how the queries are optimized and executed in a DBMS
8. to understand concurrency control and recovery in database systems\par\vspace{4pt}
\textbf{Syllabus of Courses:}\par\vspace{2pt}
Foundations\par
1. The relational model\par
2. Relational Algebra and Relational Calculus\par
3. SQL: Queries, Constraints, Triggers\par
4. Schema Refinement and Normal Forms\par
Storage and Indexing\par
1. Storing Data: Disks and Files\par
2. Tree-Structured Indexing\par
3. Hash-Based Indexing\par
Query Evaluation\par
1. External Sorting\par
2. Evaluating Relational Operators\par
3. A Typical Relational Query Optimizer\par
Transaction Management\par
1. Overview of Transaction Management\par\vspace{4pt}
\textbf{TextBooks:}\par\vspace{2pt}
1. R. Ramakrishnan and J. Gehrke, Database Management Systems, McGraw-Hill 3rd Edition,\par
2003.\par
\vspace{4pt}

\newpage
\nopagebreak
\textbf{Title of the Course:} Physical Education\\
\textbf{Number of Credits:} 1\\
\textbf{Type of the Course:} General | Practical\\
\textbf{Training hours:} 32h Practical\par\vspace{4pt}
\textbf{Course Objectives:}\par\vspace{2pt}
Physical Education classes are designed to practice and develop skills in activities that will help
you maintain fitness throughout your life. During the year, you will set and monitor personal
fitness goals that will aid in the development of a personal fitness plan.\par\vspace{4pt}

\newpage
\textbf{Title of the Course:} Compiler Design and Programming Languages\\
\textbf{Number of Credits:} 3\\
\textbf{Type of the Course:} Specialized | Theoretical\\
\textbf{Pre-requisite:} Machines and Language Theory\\
\textbf{Training hours:} 48h Theoretical\par\vspace{4pt}
\textbf{Course Objectives:}\par\vspace{2pt}
Designing and building compilers is one of the basic concepts of computer science. Although there
is little variety in the way compilers are made, a wide variety of languages and machines can be
used to make interpreters and translators. In this lesson, the subject of building compilers is
introduced by describing the main components of a compiler, their tasks, and their relationship.
After an introductory introduction to the components of a compiler and the types of grammar, the
various stages of translation, such as lexical, syntactic, and semantic analysis, and code generation
and payment, are described. Students who successfully pass this course will have a good insight
into the following: Familiarity with compiler components and various techniques for their
implementation, understanding of the implementation of programming language commands,
gaining skills in producing optimal programs and fixing programming errors, familiarity and the
use of automated tools in compiler production.\par\vspace{4pt}
\textbf{Syllabus of Courses:}\par\vspace{2pt}
An introduction to the design and implementation of programming language translators,\par
Theoretical aspects of language design and translation are discussed, and practically demonstrated\par
by developing a working compiler.\par\vspace{4pt}
\textbf{TextBooks:}\par\vspace{2pt}
1. A. Aho, M. Lam, R. Sethi, and J. Ullman, Compilers: Principles, Techniques, and Tools (2nd\par
Edition), Prentice Hall, 2006.\par
\vspace{4pt}

\newpage
\vspace*{\fill}
\begin{center}
  {\fontsize{28pt}{34pt}\selectfont \textbf{SUMMER SEMESTER}}\par\vspace{12pt}
  {\fontsize{13pt}{17pt}\selectfont Academic Year 2025–2026}
\end{center}
\vspace*{\fill}

\newpage
\nopagebreak
\textbf{Title of the Course:} Industrial Training\\
\textbf{Number of Credits:} 3\\
\textbf{Type of the Course:} Training | Practical\\
\textbf{Training hours:} 240h Practical\par\vspace{4pt}
\textbf{Course Objectives:}\par\vspace{2pt}
Students must practically get involved in industry and market by working totally specific hours in
an appropriate and defined company in either governmental or private sectors as an intern.\par\vspace{4pt}

\newpage
\vspace*{\fill}
\begin{center}
  {\fontsize{28pt}{34pt}\selectfont \textbf{SEMESTER 7}}\par\vspace{12pt}
  {\fontsize{13pt}{17pt}\selectfont Academic Year 2026–2027 (1st Semester)}
\end{center}
\vspace*{\fill}

\newpage
\textbf{Title of the Course:} Software Engineering\\
\textbf{Number of Credits:} 3\\
\textbf{Type of the Course:} Elective | Theoretical\\
\textbf{Training hours:} 48h Theoretical\par\vspace{4pt}
\textbf{Course Objectives:}\par\vspace{2pt}
The purpose of this course is to study the techniques needed to produce valuable, high-quality
high-productivity software. In this lesson, first, the differences between a product that is produced
by the engineering method and a product that is produced by the artistic methods are expressed.
Then the expectations that an engineering product has to meet are described. Afterward, we will
have an overview of scientific achievements in this field at all stages of software products with an
emphasis on engineering methods of production, including modeling, measurability, evaluability,
validation, and verification. Given the fact that in previous courses students have been less familiar
with the formal descriptions of requirements, measurements, estimates, and tests, this course
focuses more on these subjects. In the end, support activities such as project management,
scheduling, risk management, configuration management, and quality assurance, with emphasis
on their impact on software products, are reviewed. Students who successfully pass this lesson will
have a good insight into
the following:
1. Applying engineering methods to create a quality product
2. Follow the well-known processes of software engineering
3. Application of Software Architectural Design Methods
4. Software Testing at Different Levels\par\vspace{4pt}
\textbf{Syllabus of Courses:}\par\vspace{2pt}
1. An Introduction to Software Engineering\par
2. Software Development Processes and Models – Iterative Development\par
3. An overview of software analysis\par
4. Software design: design principles, patterns, refactoring\par
5. Software Architecture: Design, Documentation, and Evaluation\par
6. Software testing\par
7. Software Quality Management\par
8. Cost and time estimation\par
9. Project Management - Human Resources Management - Risk Management\par
10. Life Cycle Management - Change Management – Configuration Management\par
11. Agile methods\par\vspace{4pt}
\textbf{TextBooks:}\par\vspace{2pt}
1. Roger S. Pressman and Bruce R. Maxim, Software Engineering: A Practitioner’s Approach. 8th\par
ed., McGraw-Hill, 2014.\par
\vspace{4pt}

\newpage
\textbf{Title of the Course:} Database Laboratory\\
\textbf{Number of Credits:} 1\\
\textbf{Type of the Course:} Elective | Practical\\
\textbf{Pre-requisite:} Database Design\\
\textbf{Training hours:} 32h Practical\par\vspace{4pt}
\textbf{Course Objectives:}\par\vspace{2pt}
The course aims to cover the fundamentals of databases as seen from the perspective of application
writers. The course covers schema design techniques, SQL, data warehouses, and some aspects of
the NoSQL movement. Also, students will explain the client/server model, and describe the key
components used to implement internet database environments and perform basic database
administration tasks.\par\vspace{4pt}

\newpage
\textbf{Title of the Course:} Information Technology Project Management\\
\textbf{Number of Credits:} 3\\
\textbf{Type of the Course:} Specialized | Theoretical\\
\textbf{Training hours:} 48h Theoretical\par\vspace{4pt}
\textbf{Course Objectives:}\par\vspace{2pt}
This course provides participants with the necessary knowledge, tools, and techniques needed to
manage and control projects. It deals with all the processes necessary to plan, manage, execute,
and control projects.\par\vspace{4pt}

\newpage
\textbf{Title of the Course:} Computer Networks Security\\
\textbf{Number of Credits:} 3\\
\textbf{Type of the Course:} Elective | Theoretical\\
\textbf{Pre-requisite:} Computer Networks\\
\textbf{Training hours:} 48h Theoretical\par\vspace{4pt}
\textbf{Course Objectives:}\par\vspace{2pt}
Computer networks security is a course designed to teach students about the fundamental concepts
and techniques used in securing computer networks. With the increasing reliance on computer
networks in modern society, network security has become a crucial aspect of information
technology. This course provides an overview of the principles, practices, and technologies used
in computer networks security. The course covers a variety of topics, including cryptography,
authentication, access control, firewalls, intrusion detection and prevention, virtual private
networks (VPNs), and secure network design. Students will learn about the types of vulnerabilities
and attacks that can affect computer networks and the ways to prevent and respond to such attacks.
The goals of this course are:
1. Learning the fundamentals of cryptography
2. Learning the security protocols used in computer networks
3. Learning how security protocols are broken and hacked
Upon successful completion of the course, students will be able to
1. Do basic encryption/decryptions
2. Design secure protocols
3. Analyze the security of protocols
4. Use existing standards to secure the computer networks\par\vspace{4pt}
\textbf{Syllabus of Courses:}\par\vspace{2pt}
1- Fundamentals of Cryptography\par
2- Fundamentals of Computer Networks\par
3- Key Management\par
4- Authentication\par
5- Internet Security\par
\vspace{2pt}\textbf{- Web}\par
\vspace{2pt}\textbf{- Firewall}\par
\vspace{2pt}\textbf{- Email}\par
\vspace{2pt}\textbf{- IDS}\par
\vspace{2pt}\textbf{- VPN / IPSec}\par
\vspace{2pt}\textbf{- DNS Security}\par
6- Wireless Security\par
7- Access Control\par
8- Malware and Attacks\par\vspace{4pt}
\textbf{TextBooks:}\par\vspace{2pt}
1. “Network Security Essentials: Applications and Standards”, William Stallings, 6th ed., Pearson,\par
2017\par
\vspace{4pt}

\newpage
\textbf{Title of the Course:} Introduction to Software Testing\\
\textbf{Number of Credits:} 3\\
\textbf{Type of the Course:} Elective | Theoretical\\
\textbf{Pre-requisite:} Computer Networks\\
\textbf{Training hours:} 48h Theoretical\par\vspace{4pt}
\textbf{Course Objectives:}\par\vspace{2pt}
The purpose of this course is to acquaint students with test methods and especially test data
production methods. Also in this course, the necessary tools for the software test are introduced.\par\vspace{4pt}
\textbf{Syllabus of Courses:}\par\vspace{2pt}
Introduction to testing methods, Coverage criteria, Graph-based criteria, Criteria based on logic,\par
Criteria are based on the separation of the input space, Syntax-based criteria, Practical\par
considerations of software testing, Test tools and software setup for testing.\par\vspace{4pt}
\textbf{TextBooks:}\par\vspace{2pt}
1. P. Ammann, J. Offutt, Introduction to software testing, Cambridge University Press, 2016.\par
2. P. C. Jorgensen, Software Testing: A Craftsman's Approach, 2nd Edition, CRC Press, 2002.\par
\vspace{4pt}

\newpage
\nopagebreak
\textbf{Title of the Course:} Thematic Interpretation of the Quran\\
\textbf{Number of Credits:} 2\\
\textbf{Type of the Course:} General | Theoretical\\
\textbf{Training hours:} 32h Theoretical\par\vspace{4pt}
\textbf{Course Objectives:}\par\vspace{2pt}
Familiarity with some life-giving concepts of the Holy Quran in the form of a thematic attitude.\par\vspace{4pt}

\par\vspace{14pt}
\nopagebreak
\textbf{Title of the Course:} Sport 1\\
\textbf{Number of Credits:} 1\\
\textbf{Type of the Course:} General | Practical\\
\textbf{Training hours:} 32h Practical\par\vspace{4pt}
\textbf{Course Objectives:}\par\vspace{2pt}
This course is designed for the prospective coach, physical education, and/or recreation
professional.\par\vspace{4pt}

\end{document}
